\begin{problem}[第31届全国中学生物理竞赛复赛][Z-箍缩与安培力]{23}
Z-箍缩作为惯性约束核聚变的一种可能方式，近年来受到特别重视，其原理如图所示。图中，长 $20\mathrm{mm}$、直径为 $5\mu\mathrm{m}$ 的钨丝组成的两个共轴的圆柱面阵列，瞬间通以超强电流，钨丝阵列在安培力的作用下以极大的加速度向内运动，即所谓自箍缩效应；钨丝的巨大动量转移到处于阵列中心的直径为毫米量级的氘氚靶球上，可以使靶球压缩后达到高温高密度状态，实现核聚变。设内圈有 $N$ 根钨丝（可视为长直导线）均匀地分布在半径为 $r$ 的圆周上，通有总电流 $I_{\text{内}} = 2\times 10^{7}\mathrm{A}$；外圈有 $M$ 根钨丝，均匀地分布在半径为 $R$ 的圆周上，每根钨丝所通过的电流同内圈钨丝。已知通有电流 $i$ 的长直导线在距其 $r$ 处产生的磁感应强度大小为 $k_m \frac{i}{r}$，式中比例常量 $k_m = 2 \times 10^{-7}\mathrm{T}\cdot\mathrm{m}/\mathrm{A} = 2 \times 10^{-7}\mathrm{N}/\mathrm{A}^{2}$。
    \begin{subproblem}
        若不考虑外圈钨丝，计算内圈某一根通电钨丝中间长为 $\Delta L$ 的一小段钨丝所受到的安培力；
    \end{subproblem}
    \begin{subproblem}
        若不考虑外圈钨丝，内圈钨丝阵列熔化后形成了圆柱面，且箍缩为半径 $r = 0.25\mathrm{cm}$ 的圆柱面时，求柱面上单位面积所受到的安培力，这相当于多少个大气压？
    \end{subproblem}
    \begin{subproblem}
        证明沿柱轴方向通有均匀电流的长圆柱面，圆柱面内磁场为零，即通有均匀电流外圈钨丝的存在不改变前述两小题的结果；
    \end{subproblem}
    \begin{subproblem}
        当 $N \gg 1$ 时，则通有均匀电流的内圈钨丝在外圈钨丝处的磁感应强度大小为 $k_m \frac{I_{\text{内}}}{R}$，若要求外圈钨丝柱面每单位面积所受到的安培力大于内圈钨丝柱面每单位面积所受到的安培力，求外圈钨丝圆柱面的半径 $R$ 应满足的条件；
    \end{subproblem}
    \begin{subproblem}
        由安培环路定理可得沿柱轴方向通有均匀电流的长圆柱面外的磁场等于该圆柱面上所有电流移至圆柱轴后产生的磁场，请用其他方法证明此结论。
    \end{subproblem}
\end{problem}